\documentclass{report}
\usepackage{amsmath,amssymb}
\setlength{\parindent}{0mm}
\setlength{\parskip}{1em}
\begin{document}
\begin{center}
\rule{6in}{1pt} \
{\large Fleurianne Bertrand \\
{\bf Least-Squares mixed finite elements in relation to mixed finite elements for elasticity}}

Universit\"at Duisburg-Essen \\ Fakult�t f\"ur Mathematik \\ Thea-Leymann-Stra�e 9 \\ D-45127 Essen
\\
{\tt fleurianne.bertrand@uni-due.de}\\
Z. Cai\\
Eun Y. Park\end{center}

The related physical equations of linear elasticity are the equilibrium
equation and the constitutive equation, which expresses a relation
between the stress and strain tensors. This is a first-order partial
differential system such that a least squares method based on a
stress-displacement formulation can be used whose corresponding finite
element approximation does not preserve the symmetry of the
stress\footnote{Z. Cai, G. Starke. {\it Least squares methods for linear
elasticity.} SIAM J. Numer. Anal. 42 (2004): 826--842.}.

In this talk, a new method is investigated by introducing the vorticity
and applying the $L^2$ norm least squares principle to the
stress-displacement-vorticity system. The question of ellipticity due to
the fact that all three variables
are present in one equation is discussed. Further, the supercloseness of
the least squares approximation to the standard mixed finite element
approximations arising from the Hellinger-Reissner principle with reduced
symmetry\footnote{D. Boffi, F. Brezzi, and M. Fortin. {\it Mixed Finite
Element Methods and Applications.} Springer-Verlag, Heidelberg, 2013.
[Chp. 9]}, is studied. This implies that the favourable conservation
properties of the dual-based mixed methods and the inherent error control
of the least squares method are combined.

Additionally, a closer look will be taken at the error that appears using
this formulation on domains with curved boundaries approximated by a
triangulation \footnote{F. Bertrand, S. M\"unzenmaier, and G. Starke.
{\it First-Order System Least Squares on Curved Boundaries: Lowest-Order
Raviart-Thomas Elements.} SIAM J. Numer. Anal. 52.2 (2014): 880-894.}. In
the higher-order case, parametric Raviart-Thomas finite elements are
employed to this end.

Finally, it is shown that an optimal order of convergence is achieved and
illustrated numerically on a test example.


\end{document}
